\begin{figure}[htbp]
\centering
\begin{tikzpicture}
\begin{axis}[
    labbook standard,
    xmode=log,
    ymode=log,
    log basis x={2},
    log basis y={10},
    xlabel={Supercell Size $n_c$ (Sublattices)},
    ylabel={Cost per Map Evaluation},
    legend pos=north west,
    xmin=1.5, xmax=45,
    ymin=0.5, ymax=600
]
    \addplot[labbookteal, thick, mark=*, mark size=2.5pt]
        table[x index=0, y index=7] {../../data/processed/benchmarks/m1_7/bloch_supercell_scaling.dat};
    \addlegendentry{FLOPs per Eval [MFLOP]}

    \addplot[labbookamber, thick, mark=square*, mark size=2pt]
        table[x index=0, y index=6] {../../data/processed/benchmarks/m1_7/bloch_supercell_scaling.dat};
    \addlegendentry{Wall Time per Eval [ms]}

    % Theoretical quadratic scaling guide: O(n_c^2)
    \addplot[labbookslate, dashed, thick, domain=2:32]
        {1.573 * (x/2)^2};
    \addlegendentry{Theoretical $\mathcal{O}(n_c^2)$ Folding Law}
\end{axis}
\end{tikzpicture}
\caption{Computational scaling of the self-consistent mean-field map evaluation across generalized honeycomb supercells $n_c \in \{2, 4, 8, 16, 32\}$ under equivalent thermodynamic sampling ($N_k \cdot n_c = 1152$). The measured computational work tracks the theoretical $\mathcal{O}(n_c^2)$ folded complexity law over two orders of magnitude, with intensive ground-state energy matching to within $5.55 \times 10^{-16}$.}
\label{fig:m1_7_supercell_scaling}
\end{figure}
